\section{Threats to Validity}

\paragraph{Seed-program representativeness.}
Dataset~A may still over-represent tutorial-scale Dafny programs, particular annotation idioms, or proof styles that are easier to package into a controlled benchmark than large industrial developments. This matters primarily for generality rather than for internal causal interpretation: a benchmark can still cleanly measure brittleness within its source regimes while remaining incomplete as a census of all Dafny practice. The paper mitigates that risk by stratifying the seed corpus and by explicitly distinguishing educational examples from benchmark-style and proof-engineering-heavy programs, but this does not eliminate representativeness concerns. Any empirical pattern reported in the paper should therefore be read as evidence about the benchmarked source regimes rather than as a definitive statement about all verified Dafny code.

\paragraph{Refactoring realism.}
The benchmark uses controlled semantics-preserving refactorings instead of unconstrained maintenance history, and that choice inevitably reduces ecological realism. Real repository evolution mixes structural edits with partial rewrites, requirement changes, cleanup work, and toolchain drift in ways that are not captured by a controlled refactoring benchmark. This threat matters chiefly for external realism, not for the paper's central causal question. The paper makes this tradeoff deliberately: nearby-version pairs remain interpretable, reproducible, and suitable for clean breakage attribution, even though the benchmark does not attempt to cover the full messiness of real software evolution.

\paragraph{Semantic-preservation assumptions.}
Each transformation in Dataset~A is intended to preserve behavior, but that intention does not rule out hidden preconditions, transformation bugs, or tool-specific corner cases. A refactoring that is semantics preserving in an informal sense may still interact with Dafny parsing, name resolution, or proof-relevant structure in unexpected ways. This threat bears directly on causal interpretation because it can make an observed failure look like proof brittleness when it is partly a benchmark-construction artifact. The paper reduces this threat through seed validation, transformation preconditions, post-transformation verification replay, and exclusion of malformed cases. Even so, some residual risk remains that a subset of observed failures reflect benchmark-construction artifacts rather than pure proof brittleness.

\paragraph{Baseline sensitivity.}
The recovery outcomes for B2 and B3 depend on fixed implementation choices. For B2, rename heuristics, diff granularity, and transfer policy influence how much nearby-version continuity is recovered. For B3, prompt wording, retry limits, timeout budgets, and verifier-feedback handling influence the observed recovery rate. This is an interpretation threat rather than a benchmark-definition threat: the paper can still establish a residual repair gap, but that gap is conditional on one fixed baseline protocol. The paper addresses this by fixing a single protocol and reporting results under that protocol rather than informally tuning per instance. Even so, the measured residual repair gap should be understood as the gap that remains after these particular baseline instantiations, not as an absolute upper bound on all possible lightweight repair strategies.

\paragraph{Verifier and toolchain sensitivity.}
Some broken or unrecovered cases may reflect sensitivity of the verifier or underlying solver rather than purely source-level proof-structural difficulty. The same transformed program may behave differently under another Dafny release, solver version, or execution environment. This threat weakens causal precision for some residual failures, especially those classified as unexplained or tool-sensitive. The paper records the toolchain and environment explicitly, but it does not attempt to isolate tool-version instability as a primary benchmark axis. Accordingly, it can establish that nearby-version proof breakage is present under a fixed and documented environment, while leaving a fuller account of environment-sensitive variation to later benchmark extensions.

\paragraph{External validity.}
The conclusions of this paper should not be generalized too aggressively. Dafny is an auto-active verifier with source-visible proof structure, and the maintenance phenomena observed here may not transfer directly to interactive theorem provers, other SMT-backed verifiers, or verification ecosystems with different proof abstractions. Likewise, controlled semantics-preserving refactorings are only one slice of software evolution, and B1--B3 cover only a narrow part of the overall repair design space. The main claim is therefore intentionally modest: Dataset~A provides a reproducible maintenance-centred benchmark that makes nearby-version proof breakage measurable and exposes a residual repair gap under simple baselines.
